\section{Differences from the Original It\^o-Map Construction}\label{app:pan-differences}
We inherit from \citet{pan2026itomaps} the central definitions: the pathwise transition $\Phi_{s,t}$, its average-drift parametrization with an explicitly added martingale increment, the semigroup identity, and the velocity form of progressive self-distillation with cancelling noise. The changes below concern the process being mapped, the Brownian representation, and how the objective is trained and anchored. Table~\ref{tab:pan-differences} summarizes them. We describe the original paper only as far as the method and experiments it reports.

\paragraph{Task and conditioning.} \citet{pan2026itomaps} learn unconditional, or class-conditional, generative maps from Gaussian noise to data. They use the maps for one-shot samples of $p_{1\mid t}$ and for inference-time reward steering. We learn conditional forecast kernels $p(y_{n+1}\mid h_n)$. Every map is conditioned on a history $h$, and the generative process starts from the point mass $c(h)$ rather than from a Gaussian base. The maps are applied autoregressively over physical leads with fresh randomness per lead (Algorithm~\ref{alg:forecast}). Composition and identifiability are statements for fixed $h$ only.

\paragraph{Interpolant, SDE, and diagonal target.} The original construction uses the linear interpolant $(1-t)x_0+tx_1$ with independent Gaussian $x_0$. It maps the marginal-preserving SDE with drift $b_t+\frac{\sigma_t^2}{2}\nabla\log p_t$ and $\sigma_t=\sqrt{2(1-t)}$. Tweedie's identity then gives the diagonal target $X_1-2X_0$. That target relies on a Gaussian base independent of the data, which does not hold when $x_0=c(h)$. We instead map the forecasting SDE \eqref{eq:sde} of \citet{chen2024forecasting}, with interpolant \eqref{eq:interpolant}, schedules $\alpha_t=1-t$, $\beta_t=t^2$, and $\gamma_t=\varepsilon(1-t)$ ($\varepsilon=1$ in all experiments). Its drift $b_t$ already yields the correct marginals with diffusion $\gamma_t$, so no score correction appears. The diagonal target is $R_t$ in \eqref{eq:target}, and the diffusion is linear rather than square-root in $1-t$.

\paragraph{Brownian representation.} The original paper defaults to global Karhunen--Lo\`eve coefficients of the unweighted path on $[0,1]$, with $K=5$ per dimension, independent of $(s,t)$. It also feeds $M_t-M_s$ as an extra input. Local dyadic features of $M$ on $[s,t]$ are offered as an alternative. We use only local features: $K=4$ consecutive increments of $M$ on equal subintervals of $[s,t]$, without variance normalization. Their sum carries $M_t-M_s$. For $K$ a power of two they span the same grid values as a dyadic representation truncated at depth $\log_2K$, but with far fewer features than the original image experiments. Paths are generated lazily in the variance clock $q$ by exact Gaussian bridges and cached per example (Section~\ref{sec:map-path}). Arbitrary interval queries therefore share one realization exactly, and no fixed simulation grid or quadrature of $\sigma$ is involved. The original paper does not state how local features are coupled across the parent and child intervals of the PSD loss. Our cached bridge makes this coupling explicit, and Proposition~\ref{prop:compression} analyses its consequences.

\paragraph{Architecture and initialization.} The original maps are SiT transformers (or MLPs), trained from scratch or alongside a teacher, with Brownian features fused into the patch tokens. We wrap each experiment's existing single-time U-Net drift network. A zero-initialized interval embedding and a zero-initialized $1\times1$ path convolution are added, and all backbone weights are copied from the teacher. At initialization the student therefore equals the teacher drift, $G_\theta(s,t,x,h;\Psi)=b_{\mathrm T}(x,s;h)$.

\paragraph{Off-diagonal objective.} The original method's main objective is Lagrangian self-distillation, the residual $G+(t-s)\partial_tG-\sg(G_{t,t}(\Phi_{s,t}))$, with PSD as an alternative. We use only the composition (PSD) form, which avoids differentiating the compressed features in $t$ (Appendix~\ref{app:numerical-teacher}). Our PSD also differs in its details (Appendix~\ref{app:training-modes}):
\begin{itemize}
\item The stop-gradient target uses an EMA copy with decay $0.95$, not the current parameters as in their Algorithm~3.
\item Interval lengths follow a curriculum from $0.125$ to one, rather than $(s,u,t)$ uniform on $[0,1]$ or logit-normal times.
\item The split $u$ is restricted to the middle half of $[s,t]$.
\item An endpoint atom at $(0,\ell_{\max})$ is added, gaps below $0.01\ell_{\max}$ are excluded, and the consistency weight is ramped.
\end{itemize}

\paragraph{Role of the teacher.} The original teacher-guided variant uses a pretrained deterministic flow map, Decoupled MeanFlow. Its velocity is converted into the stochastic diagonal through $(1+t)v_{t,t}-x$, and it supervises the diagonal and the Lagrangian target. The original paper notes that PSD may use the same teacher diagonal. Our teacher is the stochastic-interpolant drift of the same forecasting SDE, so it supplies $b_{\mathrm T}$ directly without conversion. Our stage one is the PSD-with-teacher-diagonal variant. In addition, our teacher initializes the student and selects checkpoints through same-path teacher rollouts.

\paragraph{Simulation-based refinement.} Neither of the original objectives simulates trajectories. Our second stage regresses on the same-path numerical teacher average drift \eqref{eq:tr-target}: 32 Euler drift steps with exact martingale increments. It has no counterpart in the original method. It gives up a fully simulation-free training claim in exchange for direct finite-time supervision.

\paragraph{Deterministic comparator.} The original paper compares against flow-map and Meta Flow Map steering. Our deterministic baseline is a history-conditioned flow map from a Gaussian base. It is trained with the same two-stage procedure and interface, with the Brownian features and martingale term removed (Section~\ref{sec:flow-map-baseline}).

\paragraph{Theoretical claims.} The original identification results (their Theorem~2 and Corollary~1) assume the full Brownian path, full-support sampling of $(s,t,x,W)$, and an exact diagonal. Our training violates each of these. The features are finite, intervals are restricted and curriculum-sampled, the target is an EMA, and the diagonal is a learned teacher. We therefore claim only the ideal identities for fixed $h$. In place of identifiability, we add the elementary same-path coupling bound and the compression decomposition \eqref{eq:compression}, with diagnostics quantifying the latter.

\paragraph{Evaluation.} The original experiments evaluate same-path fidelity (Gaussian mixtures, MNIST) and reward steering (Gaussian mixtures, MNIST, ImageNet-256), without sample-quality metrics such as FID. We evaluate held-out probabilistic forecasts of stochastically forced Navier--Stokes: fair CRPS and energy score, spread--skill and calibration diagnostics, enstrophy spectra, and long autoregressive rollouts. Measured sampling latency is reported alongside quality against multi-step stochastic-interpolant and flow-matching samplers. Inference-time steering is not studied.

\begin{table*}[t]
\caption{Differences between the original It\^o-map construction \citep{pan2026itomaps} and the conditional forecasting maps studied here.}
\label{tab:pan-differences}
\centering\small
\begin{tabular}{@{}p{0.17\textwidth}p{0.385\textwidth}p{0.385\textwidth}@{}}
\toprule
 & \citet{pan2026itomaps} & This work\\
\midrule
Target law & $p_{\mathrm{data}}$ (optionally class-conditional) & $p(y_{n+1}\mid h_n)$, history-conditioned, autoregressive\\
Initial state & Gaussian $x_0$, independent of $x_1$ & Point mass $c(h)$\\
Process & $dX=(b+\tfrac{\sigma^2}{2}\nabla\log p)\,dt+\sigma\,dW$, $\sigma_t=\sqrt{2(1-t)}$ & Forecasting SDE \eqref{eq:sde}, $\gamma_t=\varepsilon(1-t)$\\
Diagonal target & $X_1-2X_0$ & $R_t=-c(h)+2ty-\varepsilon\sqrt t z$, or $b_{\mathrm T}$\\
Path features & Global KL ($K{=}5$) plus $M_t-M_s$; or local dyadic & $K{=}4$ local increments of $M$; cached variance-clock bridges\\
Off-diagonal loss & Lagrangian (main) or PSD & PSD only\\
PSD target & Current parameters, stop-gradient & EMA ($0.95$), stop-gradient\\
Time sampling & Uniform or logit-normal $(s,t)$ & Interval curriculum, middle-half split, endpoint atom\\
Teacher & Optional deterministic DMF, converted to the stochastic diagonal & Same-SDE drift: diagonal target, initialization, validation\\
Simulation & None & Stage two: 32-step same-path teacher rollouts\\
Network & SiT or MLP & Teacher U-Net with zero-initialized interval and path branches\\
Theory & Identifiability under full path, full support, exact diagonal & Ideal identities only; coupling and compression bounds\\
Evaluation & Same-path fidelity, reward steering & Navier--Stokes forecast CRPS, calibration, spectra, long rollouts, latency\\
\bottomrule
\end{tabular}
\end{table*}
